\chapter{Conclusion and Future Work}
\label{ch:conc}

\section{Conclusion}

We built a complete post-quantum security framework for connected vehicles,
from registration with a Trusted Authority through a mutually authenticated
ML-KEM/ML-DSA handshake and encrypted records to AGS-PBFT consensus over a
verifiable ledger, and ran it as a distributed system, in an 802.11p
simulation with mobility, and under attack. Implementing the base paper
exposed defects that prevent it from working as published, the most serious
being a consensus layer that signs with a KEM; implementing our own Review~1
design exposed the lack of client authentication. Both were corrected and the
corrections tested.

The measurements answer the project's question. Post-quantum cryptography
fits the vehicular latency budget per message, at every NIST security level,
once it is implemented natively: a beacon inside a session is a
\V{deck-beacon-aead-bytes}-byte AES-GCM record whatever the suite, and a full
handshake takes about \V{deck-hs-ms}\,ms of computation. The binding
constraint is size. A post-quantum handshake is about twenty times larger
than a classical one, and on a shared 6\,Mb/s channel that is what fails as
density grows: at 100 vehicles in our scenario DECK completes
\V{ns3-success-deck-100}\,\% of handshakes, against
\V{ns3-success-classic-100}\,\% classically. Protocol design recovers much of
the loss. Certificate caching halves the handshake and raises completion to
\V{ns3-success-deck-cache-100}\,\%, and exponential backoff after a failed
handshake prevents congestion collapse. Safety warnings do not wait
for consensus at all, and AGS-PBFT saves messages, and CPU at light load, but not time.

\section{Limitations}

\begin{itemize}
\item \textbf{One machine.} The distributed system runs on localhost, so its
latencies include TCP over loopback, not radio. Radio effects come only from
ns-3.
\item \textbf{ns-3 replays measured costs.} The simulation uses measured
message sizes and processing-time distributions rather than executing the
cryptography.
\item \textbf{No WAVE module.} The ns-3 Python package provides the 802.11p
PHY and MAC but no WAVE/OCB module, so the ad-hoc MAC is used. Ad-hoc mode
also has no association, so behaviour is close.
\item \textbf{Certificate caching relies on infrastructure we model.} The
edge directory is the TA registry in the distributed system and one backhaul
round trip in ns-3. A vehicle's first handshake still sends its certificate in
full.
\item \textbf{Linkability.} Certificates and certificate references identify a
vehicle across handshakes; pseudonym certificates would be needed for
location privacy.
\item \textbf{V2V through RSUs is not end-to-end.} Relaying RSUs see the
plaintext, as in the original design.
\item \textbf{No formal verification} of the handshake.
\end{itemize}

\section{Future work}

\begin{itemize}
\item \textbf{Smaller handshakes.} KEM-based authentication in the style of
KEMTLS, which replaces the server's signature with a KEM, or certificates
signed with a scheme that has smaller signatures, would attack the remaining
size directly.
\item \textbf{Robust regrouping.} Replace AGS-PBFT's $\mu \pm \sigma$ rule
with the median and median absolute deviation, so that one outlier cannot
shield a silent node.
\item \textbf{Formal analysis} of the handshake, including certificate
caching, with a protocol verifier such as ProVerif or Tamarin.
\item \textbf{Pseudonymous certificates} compatible with caching.
\item \textbf{Realistic mobility} from a traffic simulator on a real road
network, and evaluation on vehicular hardware.
\end{itemize}
